\documentclass[11pt,a4paper]{article}

% ── Codificación, fuentes y lenguaje ───────────────────────────────────────────
\usepackage[utf8]{inputenc}
\usepackage[T1]{fontenc}
\usepackage[default,scale=0.95]{sourcesanspro}
\renewcommand{\familydefault}{\sfdefault}
\usepackage[spanish,es-noshorthands,es-tabla]{babel}

% ── Geometría y espaciado ─────────────────────────────────────────────────────
\usepackage[top=2.2cm, bottom=2.5cm, left=2.2cm, right=2.2cm, headheight=15pt]{geometry}
\usepackage{setspace}
\setstretch{1.18}
\usepackage{parskip}

% ── Matemáticas y unidades ────────────────────────────────────────────────────
\usepackage{amsmath}
\usepackage{amssymb}
\usepackage{siunitx}

% ── Circuitos ─────────────────────────────────────────────────────────────────
\usepackage[european, straightvoltages]{circuitikz}

% ── Tablas ────────────────────────────────────────────────────────────────────
\usepackage{booktabs}
\usepackage{tabularx}
\usepackage{array}
\usepackage{longtable}
\usepackage{needspace}
\usepackage{colortbl}
\usepackage{enumitem}

% ── Colores, cajas e iconos ───────────────────────────────────────────────────
\usepackage[dvipsnames]{xcolor}
\usepackage{tcolorbox}
\tcbuselibrary{skins, breakable}
\usepackage{fontawesome5}
\usepackage{graphicx}

% ── Listados de código (dark-mode infográfico) ────────────────────────────────
%   Estilo base: fondo oscuro con números de línea. Los bloques lstlisting se
%   envuelven automáticamente en una caja tcolorbox redondeada.
\usepackage{listings}

% ── Secciones con barra de color (infografía) ────────────────────────────────
\usepackage{titlesec}
\titleformat{\section}
  {\normalfont\LARGE\bfseries\color{azulNoche}}
  {\colorbox{cyanNeon}{\color{fondoOscuro}\thesection}}{0.7em}{}
  [\vspace{2pt}\color{cyanNeon}\rule{\linewidth}{1.8pt}]
\titlespacing*{\section}{0pt}{24pt}{10pt}
\titleformat{\subsection}
  {\normalfont\large\bfseries\color{azulNoche}}
  {\textcolor{cyanNeon}{\thesubsection}}{0.7em}{}
  [\vspace{1pt}\color{grisLinea}\rule{\linewidth}{0.6pt}]
\titlespacing*{\subsection}{0pt}{14pt}{6pt}
\titleformat{\subsubsection}
  {\normalfont\normalsize\bfseries\color{azulNoche}}
  {\textcolor{cyanNeon}{\thesubsubsection}}{0.7em}{}
\titlespacing*{\subsubsection}{0pt}{10pt}{4pt}

% ── Cabeceras y pies de página ────────────────────────────────────────────────
\usepackage{fancyhdr}
\usepackage{lastpage}
\pagestyle{fancy}
\fancyhf{}
\renewcommand{\headrulewidth}{0pt}
\renewcommand{\footrulewidth}{0pt}
\fancyfoot[C]{%
  \begin{minipage}{\linewidth}
    \vspace{2pt}
    \color{cyanNeon}\rule{\linewidth}{1.2pt}\\[2pt]
    \centering\color{grisTexto}\footnotesize
    Página \thepage\ de \pageref{LastPage}
  \end{minipage}%
}

% ── Hiperenlaces ──────────────────────────────────────────────────────────────
\usepackage[colorlinks=true, linkcolor=azulNoche, urlcolor=cyanNeon]{hyperref}
\sloppy
\emergencystretch 3em

% ── Paleta de colores — tecnológico dark-mode ──────────────────────────────────
\definecolor{fondoOscuro}{HTML}{0D1B2A}%             Dark navy — fondo banda título
\definecolor{verdeTurquesa}{HTML}{004D40}%           Verde turquesa — bandas de marca
\definecolor{azulNoche}{HTML}{1B3A6B}%               Midnight blue — secciones, marcos
\definecolor{azulMedio}{HTML}{2A5298}%               Azul medio — variante de acento
\definecolor{cyanNeon}{HTML}{00C8E0}%                Cyan eléctrico — acentos primarios
\definecolor{verdeSignal}{HTML}{00B887}%             Verde señal — fortalezas, éxito
\definecolor{naranjaVivo}{HTML}{FF7043}%             Naranja vivo — mejoras, advertencias
\definecolor{rojoAlerta}{HTML}{E53935}%              Rojo alerta — errores
\definecolor{grisPapel}{HTML}{F4F6FA}%               Fondo tarjetas claras
\definecolor{grisLinea}{HTML}{C8D0E0}%               Bordes sutiles
\definecolor{grisTexto}{HTML}{6B7A99}%               Texto secundario
\definecolor{amarilloNota}{HTML}{FFB300}%            Notas / advertencias doradas
\definecolor{codigoFondo}{HTML}{1E2D3D}%             Fondo del bloque de código
\definecolor{codigoComentario}{HTML}{7A8FA6}%        Comentarios
\definecolor{codigoCadena}{HTML}{FF9D5C}%            Cadenas / strings
\definecolor{codigoKeyword}{HTML}{00C8E0}%           Palabras clave
\definecolor{codigoNumero}{HTML}{00E5A0}%            Números

% ── Banda de título: portada infográfica ──────────────────────────────────────
%   Diseño en 2 capas: banda de color (por defecto fondo oscuro) + franja de
%   acento cyan abajo. Uso: \bandaTitulo[color]{título}{subtítulo}.
%   El icono se puede cambiar con \renewcommand{\iconoBanda}{...} (FontAwesome).
\newcommand{\iconoBanda}{microchip}
\newcommand{\bandaTitulo}[3][fondoOscuro]{%
  \begin{tcolorbox}[
    enhanced,
    colback=#1, colframe=#1,
    boxrule=0pt, arc=8pt, outer arc=8pt,
    width=\linewidth,
    top=18pt, bottom=6pt, left=14pt, right=14pt,
    borderline south={3.5pt}{0pt}{cyanNeon},
  ]
    \begin{minipage}[c]{0.07\linewidth}
      \centering
      \textcolor{cyanNeon}{\fontsize{30}{30}\selectfont\faIcon{\iconoBanda}}
    \end{minipage}%
    \hspace{8pt}%
    \begin{minipage}[c]{0.88\linewidth}
      {\fontsize{20}{24}\selectfont\bfseries\color{white}#2}\par\vspace{5pt}%
      \textcolor{cyanNeon!80!white}{\small\faIcon{angle-right}~#3}%
    \end{minipage}
    \vspace{4pt}
  \end{tcolorbox}%
  \vspace{8pt}%
}

% ── Tarjetas de datos (infografía) ────────────────────────────────────────────
\tcbset{
  tarjetaDato/.style={
    enhanced, breakable,
    arc=6pt, outer arc=6pt,
    colback=grisPapel, colframe=azulNoche,
    boxrule=1pt,
    fonttitle=\bfseries\small, coltitle=white,
    attach boxed title to top left={yshift=-3mm, xshift=6mm},
    boxed title style={
      colback=azulNoche, colframe=azulNoche,
      arc=4pt, boxrule=0pt,
    },
    drop fuzzy shadow=azulNoche!30!white,
    top=8pt, bottom=8pt, left=10pt, right=10pt,
  },
  tarjetaIcono/.style={
    enhanced, breakable,
    arc=6pt, outer arc=6pt,
    colback=white, colframe=grisLinea, boxrule=0.8pt,
    fonttitle=\bfseries\small, coltitle=azulNoche,
    borderline west={3pt}{0pt}{cyanNeon},
    drop fuzzy shadow=azulNoche!25!white,
    top=6pt, bottom=6pt, left=10pt, right=10pt,
  },
  cajaTitulo/.style={
    enhanced, breakable,
    arc=6pt, outer arc=6pt,
    colback=azulNoche!10!grisPapel, colframe=azulNoche,
    boxrule=1pt,
    fonttitle=\bfseries, coltitle=white,
    attach boxed title to top left={yshift=-3mm, xshift=6mm},
    boxed title style={
      colback=azulNoche, colframe=azulNoche,
      arc=4pt, boxrule=0pt,
    },
    drop fuzzy shadow=azulNoche!25!white,
    top=8pt, bottom=8pt, left=10pt, right=10pt,
  },
  cajaContenido/.style={
    enhanced, breakable,
    arc=5pt, outer arc=5pt,
    colback=grisPapel, colframe=grisLinea,
    boxrule=0.8pt,
    fonttitle=\bfseries\small, coltitle=azulNoche,
    top=6pt, bottom=6pt, left=10pt, right=10pt,
  },
  cajaFortaleza/.style={
    enhanced, breakable,
    arc=6pt, outer arc=6pt,
    colback=verdeSignal!8!white, colframe=verdeSignal,
    boxrule=1pt,
    borderline west={4pt}{0pt}{verdeSignal},
    fonttitle=\bfseries\small, coltitle=verdeSignal!70!black,
    drop fuzzy shadow=verdeSignal!20!white,
    top=6pt, bottom=6pt, left=10pt, right=10pt,
  },
  cajaMejora/.style={
    enhanced, breakable,
    arc=6pt, outer arc=6pt,
    colback=naranjaVivo!8!white, colframe=naranjaVivo,
    boxrule=1pt,
    borderline west={4pt}{0pt}{naranjaVivo},
    fonttitle=\bfseries\small, coltitle=naranjaVivo!80!black,
    drop fuzzy shadow=naranjaVivo!20!white,
    top=6pt, bottom=6pt, left=10pt, right=10pt,
  },
  cajaRecomendacion/.style={
    enhanced, breakable,
    arc=6pt, outer arc=6pt,
    colback=cyanNeon!6!white, colframe=cyanNeon!60!azulNoche,
    boxrule=1pt,
    borderline west={4pt}{0pt}{cyanNeon},
    fonttitle=\bfseries\small, coltitle=azulNoche,
    drop fuzzy shadow=azulNoche!20!white,
    top=6pt, bottom=6pt, left=10pt, right=10pt,
  },
}

% ── Ícono + texto (encabezados de sección y elementos) ────────────────────────
\newcommand{\iconotexto}[2]{\textcolor{cyanNeon}{\faIcon{#1}}~#2}

% ── Listados de código: estilo dark + auto-envoltura ─────────────────────────
\lstdefinestyle{estiloCodigo}{
    backgroundcolor=\color{codigoFondo},
    basicstyle=\ttfamily\small\color{white},
    commentstyle=\color{codigoComentario}\itshape,
    stringstyle=\color{codigoCadena},
    keywordstyle=\color{codigoKeyword}\bfseries,
    numberstyle=\tiny\color{codigoComentario},
    numbers=left,
    numbersep=10pt,
    showstringspaces=false,
    breaklines=true,
    breakatwhitespace=false,
    tabsize=4,
    frame=none,
    captionpos=b,
    aboveskip=6pt,
    belowskip=4pt,
    xleftmargin=14pt,
    framexleftmargin=14pt,
    literate=
      {á}{{\'a}}1 {é}{{\'e}}1 {í}{{\'i}}1 {ó}{{\'o}}1 {ú}{{\'u}}1
      {Á}{{\'A}}1 {É}{{\'E}}1 {Í}{{\'I}}1 {Ó}{{\'O}}1 {Ú}{{\'U}}1
      {ñ}{{\~n}}1 {Ñ}{{\~N}}1 {ü}{{\"u}}1 {Ü}{{\"U}}1
      {¿}{{?`}}1 {¡}{{!`}}1
}
\lstdefinestyle{estiloPython}{
    style=estiloCodigo,
    language=Python,
}
\lstset{style=estiloCodigo}
\BeforeBeginEnvironment{lstlisting}{%
  \begin{tcolorbox}[
    enhanced, breakable,
    arc=6pt, outer arc=6pt,
    colback=codigoFondo, colframe=azulNoche,
    boxrule=1pt,
    drop fuzzy shadow=azulNoche!25!white,
    top=2pt, bottom=2pt, left=0pt, right=0pt,
  ]%
}
\AfterEndEnvironment{lstlisting}{\end{tcolorbox}}

% ── Cajas de contenido temáticas ──────────────────────────────────────────────
\newtcolorbox{cajaRecuerda}{
    enhanced, breakable,
    arc=6pt, outer arc=6pt,
    colback=azulNoche!10!grisPapel,
    colframe=azulNoche, boxrule=1pt,
    borderline west={4pt}{0pt}{cyanNeon},
    fonttitle=\bfseries\small\color{white},
    title={\faIcon{info-circle}~Recuerda},
    attach boxed title to top left={yshift=-3mm, xshift=6mm},
    boxed title style={colback=azulNoche, arc=4pt, boxrule=0pt},
    drop fuzzy shadow=azulNoche!25!white,
    left=10pt, right=10pt, top=8pt, bottom=8pt,
}

\newtcolorbox{cajaConcepto}{
    enhanced, breakable,
    arc=6pt, outer arc=6pt,
    colback=verdeSignal!8!white,
    colframe=verdeSignal, boxrule=1pt,
    borderline west={4pt}{0pt}{verdeSignal},
    fonttitle=\bfseries\small\color{white},
    title={\faIcon{lightbulb}~Concepto clave},
    attach boxed title to top left={yshift=-3mm, xshift=6mm},
    boxed title style={colback=verdeSignal!80!black, arc=4pt, boxrule=0pt},
    drop fuzzy shadow=verdeSignal!20!white,
    left=10pt, right=10pt, top=8pt, bottom=8pt,
}

\newtcolorbox{cajaEjemplo}{
    enhanced, breakable,
    arc=6pt, outer arc=6pt,
    colback=cyanNeon!5!white,
    colframe=cyanNeon!70!azulNoche, boxrule=1pt,
    borderline west={4pt}{0pt}{cyanNeon},
    fonttitle=\bfseries\small\color{fondoOscuro},
    title={\faIcon{code}~Ejemplo},
    attach boxed title to top left={yshift=-3mm, xshift=6mm},
    boxed title style={colback=cyanNeon, arc=4pt, boxrule=0pt},
    drop fuzzy shadow=azulNoche!20!white,
    left=10pt, right=10pt, top=8pt, bottom=8pt,
}

\newtcolorbox{cajaEjercicio}{
    enhanced, breakable,
    arc=6pt, outer arc=6pt,
    colback=azulMedio!7!white,
    colframe=azulMedio, boxrule=1pt,
    borderline west={4pt}{0pt}{azulMedio},
    fonttitle=\bfseries\small\color{white},
    title={\faIcon{pencil-alt}~Ejercicio},
    attach boxed title to top left={yshift=-3mm, xshift=6mm},
    boxed title style={colback=azulMedio, arc=4pt, boxrule=0pt},
    drop fuzzy shadow=azulNoche!20!white,
    left=10pt, right=10pt, top=8pt, bottom=8pt,
}

\newtcolorbox{cajaReto}{
    enhanced, breakable,
    arc=6pt, outer arc=6pt,
    colback=naranjaVivo!6!white,
    colframe=naranjaVivo, boxrule=1pt,
    borderline west={4pt}{0pt}{naranjaVivo},
    fonttitle=\bfseries\small\color{white},
    title={\faIcon{fire}~Reto},
    attach boxed title to top left={yshift=-3mm, xshift=6mm},
    boxed title style={colback=naranjaVivo, arc=4pt, boxrule=0pt},
    drop fuzzy shadow=naranjaVivo!20!white,
    left=10pt, right=10pt, top=8pt, bottom=8pt,
}

% ── Indicadores de dificultad ─────────────────────────────────────────────────
\newcommand{\facil}{\textcolor{verdeSignal}{\faIcon{circle}\,\textbf{Fácil}}}
\newcommand{\intermedio}{\textcolor{naranjaVivo}{\faIcon{circle}\,\textbf{Intermedio}}}
\newcommand{\dificil}{\textcolor{rojoAlerta}{\faIcon{circle}\,\textbf{Difícil}}}

% ─────────────────────────────────────────────────────────────────────────────
\renewcommand{\iconoBanda}{shield-alt}
\fancyhead[L]{\footnotesize\color{grisTexto}\faIcon{shield-alt}~CYBERSEGURIDAD}
\fancyhead[R]{\footnotesize\color{grisTexto}Semana 07}
\hypersetup{pdftitle={Ciberseguridad Aplicada},pdfauthor={Prof. César Rodríguez}
  pdfsubject={CYBERSEGURIDAD}}

\begin{document}
\bandaTitulo[fondoOscuro]{Semana 07 --- Repetir el ataque y demostrar que ya no funciona}{Grupo 3 $\cdot$ rol Rojo $\cdot$ ejemplo de trabajo}

\begin{center}\small\color{grisTexto}Prof. César Rodríguez \textbar{} ELECTRONICA \textbar{} 1 de noviembre de 2026\end{center}

\begin{tabularx}{\linewidth}{@{}>{\raggedright\arraybackslash}p{4.6cm}X@{}}
\toprule
\textbf{Curso} & CYBERSEGURIDAD --- Ciberseguridad Aplicada \\
\textbf{Semana} & 07 de 10 \\
\textbf{Fase} & 3 --- Defensa y hardening \\
\textbf{Grupo} & 3 \\
\textbf{Rol} & Rojo \\
\textbf{Entrega} & GitHub, carpeta semana\_07/ \\
\textbf{Corte} & Domingo 23:59 de la semana correspondiente; lo posterior no se corrige. \\
\bottomrule
\end{tabularx}

\needspace{5\baselineskip}
\section{\iconotexto{user-shield}{Tu rol en la práctica}}

Ejecuta la técnica ofensiva sobre el objetivo autorizado: prepara el entorno, lanza el ataque, registra comandos y payloads, y guarda cada salida con su hash.

\begin{itemize}[leftmargin=1.4em,itemsep=2pt,topsep=3pt]
  \item Confirmar el objetivo autorizado y preparar el entorno antes de la primera captura.
  \item Ejecutar la técnica de la semana en el orden de los pasos y guardar cada salida cruda con su hash.
  \item Redactar la reproducción: comando, salida esperada y salida obtenida, sin retocar la evidencia.
\end{itemize}

\needspace{5\baselineskip}
\section{\iconotexto{flag}{Qué aporta este rol al entregable}}

Evidencia cruda en semana\_NN/evidencia y el apartado de reproducción del informe.

\needspace{5\baselineskip}
\section{\iconotexto{list-ol}{Ejemplo paso a paso}}

El rol Rojo conserva los payloads de la semana 5 y los repite contra la copia parcheada. La prueba es el mismo ataque, no uno parecido.

\subsection{Guardar la evidencia del antes en la copia sin parchear}

Se reproduce cada hallazgo y se guardan peticion y respuesta. Esta evidencia es la referencia contra la que se compara el despues.

\begin{lstlisting}[style=estiloCodigo]
mkdir -p semana_07/evidencia
curl -s -b semana_07/evidencia/cookies.txt \
  'http://127.0.0.1:8081/vulnerabilities/sqli/?id=1%27+OR+%271%27%3D%271&Submit=Submit' \
  > semana_07/evidencia/sqli_antes.html
grep -c 'First name' semana_07/evidencia/sqli_antes.html \
  > semana_07/evidencia/sqli_antes_conteo.txt
\end{lstlisting}

\subsection{Reaplicar el mismo payload sobre el codigo parcheado}

Se cambia solo el destino: mismo parametro, mismo orden, misma sesion. Si el resultado ya no devuelve filas extra, la mitigacion se sostiene.

\begin{lstlisting}[style=estiloCodigo]
curl -s -b semana_07/evidencia/cookies.txt \
  'http://127.0.0.1:8081/vulnerabilities/sqli/?id=1%27+OR+%271%27%3D%271&Submit=Submit' \
  > semana_07/evidencia/sqli_despues.html
grep -c 'First name' semana_07/evidencia/sqli_despues.html \
  > semana_07/evidencia/sqli_despues_conteo.txt
diff semana_07/evidencia/sqli_antes_conteo.txt \
  semana_07/evidencia/sqli_despues_conteo.txt || true
\end{lstlisting}

\subsection{Probar una variante del mismo ataque}

Cambiar mayusculas, comentarios o codificacion es parte del mismo ataque. Si la variante funciona, el parche cubrio el ejemplo y no la causa.

\subsection{Documentar la regresion funcional}

Ademas de que el ataque falle, la funcionalidad legitima debe seguir funcionando. Se prueban los casos normales del formulario.

\needspace{8\baselineskip}\subsection{Hallazgo, payload y resultado despues}

\begin{center}
\begin{minipage}{\linewidth}
\centering
{\footnotesize
\textbf{Hallazgo, payload y resultado despues}\par\smallskip
\begin{tabular}{@{}>{\raggedright\arraybackslash}p{2.20cm}>{\raggedright\arraybackslash}p{3.60cm}>{\raggedright\arraybackslash}p{2.40cm}>{\raggedright\arraybackslash}p{3.20cm}>{\raggedright\arraybackslash}p{2.60cm}@{}}
\toprule
\textbf{Hallazgo} & \textbf{Payload original} & \textbf{Resultado} & \textbf{Variante probada} & \textbf{Conclusion} \\
\midrule
Inyeccion SQL & Comilla y OR tautologico & Sin filas extra & Mayusculas y comentario & Mitigado \\
XSS reflejado & Script en name & Texto codificado & Etiqueta img & Mitigado \\
LFI & Secuencia ../ en page & Error controlado & Codificacion URL & Mitigado \\
Formulario legitimo & Consulta normal & Funciona igual & No aplica & Sin regresion \\
\bottomrule
\end{tabular}
}
\par\smallskip{\footnotesize La columna de variante se llena solo si se probo; si no, se marca como no probada en lugar de suponer que tambien falla.}
\end{minipage}
\end{center}

\needspace{5\baselineskip}
\section{\iconotexto{box-open}{Qué guardar en el repositorio}}

Se entrega la evidencia antes y despues con el mismo payload, el resultado de la variante y la comprobacion de que la funcionalidad legitima sigue operativa.

\begin{lstlisting}[style=estiloCodigo]
semana_07/
|-- informe.md
|-- codigo/
|-- evidencia/
|   |-- sqli_antes.html
|   |-- sqli_despues.html
|   |-- sqli_antes_conteo.txt
|   `-- sqli_despues_conteo.txt
`-- README.md
\end{lstlisting}

\begin{itemize}[leftmargin=1.4em,itemsep=2pt,topsep=3pt]
  \item El payload de antes y el de despues son identicos.
  \item La variante se probo o se declaro no probada.
  \item La funcionalidad legitima se verifico despues del parche.
  \item La conclusion distingue mitigado de no verificado.
\end{itemize}

\needspace{5\baselineskip}
\section{\iconotexto{ruler-combined}{Cómo se corrige}}

\begin{tabularx}{\linewidth}{@{}X r@{}}
\toprule
Criterio & Peso\\
\midrule
Correcciones que eliminan la causa raíz, no el síntoma & 3,0 \\
Demostración antes y después con el mismo payload & 2,5 \\
Controles automatizados ejecutados e interpretados & 2,0 \\
Criterio de aceptación que previene la regresión & 1,5 \\
Sin regresiones funcionales en la aplicación & 1,0 \\
\midrule
\textbf{Total} & \textbf{10,0}\\
\bottomrule
\end{tabularx}

\begin{itemize}[leftmargin=1.4em,itemsep=2pt,topsep=3pt]
  \item Cambiar el payload para que falle no demuestra nada.
  \item Un parche que rompe la aplicacion no es una mitigacion aceptable.
\end{itemize}

\begin{cajaRecuerda}El trabajo es independiente: cada grupo usa su propio repositorio y su propio laboratorio, y la comparación de resultados es en clase. Corte: Domingo 23:59 de la semana correspondiente; lo posterior no se corrige.
\end{cajaRecuerda}
\end{document}
